\section{Step 3 Inspect the supplied leaf entries and returns}\label{rgb-s03}

Each entry copies the incoming count to \texttt{v0} before \texttt{a0}, \texttt{a1} and \texttt{a2} become working pointers and remaining count. The functions are leaves, so their incoming \texttt{ra} remains live. No saved register or native stack slot is needed by this guided design; equivalent nonleaf implementations must preserve their own return and frame correctly.

\subsection{Expected observation}
The provided exits return through \texttt{ra}. Temporary vector gates initially choose the scalar path; other unfinished regions choose their function return.

\Needspace{11\baselineskip}
\subsection{Check your understanding}
Why retain the count before the loop?

\begin{itemize}
  \item \textbf{A} To set the memory allocation size.
  \item \textbf{B} To choose a native return address.
  \item \textbf{C} The remaining count will be consumed while the public result stays original.
\end{itemize}
